% GENERATED FILE. Edit canonical lessons, then run npm run generate:books.
% canonical-source-hash: fnv1a64:e6feb8aac4395a01
% canonical-lessons: ZH-C257-yueqi, ZH-C257-dagu, ZH-C257-yanchu, ZH-C257-geshou, ZH-C257-guzhang

\chapter{Music and Shows}
\label{ch:zh-music-and-shows}
\hlchaptermodality{257}

\begin{chapteropening}
\textbf{By the end of this chapter:} \emph{I can talk about musical instruments, drumming, performing and singers, and say people clapped.}

Complete the last lesson of chapter 257: I can talk about musical instruments, drumming, performing and singers, and say people clapped.
\end{chapteropening}

\section[yuèqì]{yuèqì — musical instrument}

\label{lesson:ZH-C257-yueqi}



\begin{quote}
Before the new one: say the Mandarin for it may be that, then the Mandarin for different.
\end{quote}

\subsection*{You'll want to know: yuèqì}
\emph{yuèqì} — \textquotedblleft{}musical instrument\textquotedblright{}.

\begin{quote}
\emph{nǐ huì tán shénme yuèqì} — What instrument can you play?
\end{quote}

Its characters come later; for now it is a word for the ear.

\begin{cousinweb}
\emph{yuèqì} is built from \zh{乐} 'music' and \emph{qì} 'a vessel', both inherited from Middle Chinese, \emph{ngaewk} and \emph{khijH} (Guangyun, in Baxter's transcription), which glosses the first 'music' and the second 'a dish, a vessel'.
\end{cousinweb}

\begin{grammarlens}[title={What you've built}]
One more word for what I like.
\end{grammarlens}

\subsection*{Your turn}
\begin{itemize}
\raggedright
  \item \emph{Say it:} \emph{yuèqì}
  \item \emph{Say it:} \emph{yuèqì}, once more
  \item \emph{Say it:} \emph{nǐ huì tán shénme yuèqì} — What instrument can you play?
  \item \emph{Recall:} say the Mandarin for it may be that, then the Mandarin for different, then say \emph{yuèqì} again
\end{itemize}

\begin{tcolorbox}[breakable,colback=teal!4,colframe=teal!35!black,title={Before you move on}]
What does yuèqì mean? (\textquotedblleft{}Musical instrument\textquotedblright{}.) Say it once more.
\end{tcolorbox}
\section[dǎgǔ]{dǎgǔ — to play the drums}

\label{lesson:ZH-C257-dagu}



\begin{quote}
Before the new one: say the Mandarin for different, then the Mandarin for a musical instrument.
\end{quote}

\subsection*{You'll want to know: dǎgǔ}
\emph{dǎgǔ} — \textquotedblleft{}to play the drums\textquotedblright{}.

\begin{quote}
\emph{wǒ dìdi huì dǎgǔ} — My little brother can play the drums.
\end{quote}

Its characters come later; for now it is a word for the ear.

Two third tones meet in this word: dǎ is said with a rising tone before gǔ, and the spelling keeps both marks.

\begin{cousinweb}
\emph{dǎgǔ} is built from \zh{打} 'to strike' and \emph{gǔ} 'a drum', both inherited from Middle Chinese, \emph{taengX} and \emph{kuX} (Guangyun, in Baxter's transcription), which quotes an older dictionary for the second, 'to beat a drum'. Middle Chinese \emph{taengX} would regularly give dǎng; dǎ is irregular.
\end{cousinweb}

\begin{grammarlens}[title={What you've built}]
One more word for what I like.
\end{grammarlens}

\subsection*{Your turn}
\begin{itemize}
\raggedright
  \item \emph{Say it:} \emph{dǎgǔ}
  \item \emph{Say it:} \emph{dǎgǔ}, once more
  \item \emph{Say it:} \emph{wǒ dìdi huì dǎgǔ} — My little brother can play the drums.
  \item \emph{Recall:} say the Mandarin for different, then the Mandarin for a musical instrument, then say \emph{dǎgǔ} again
\end{itemize}

\begin{tcolorbox}[breakable,colback=teal!4,colframe=teal!35!black,title={Before you move on}]
What does dǎgǔ mean? (\textquotedblleft{}To play the drums\textquotedblright{}.) Say it once more.
\end{tcolorbox}
\section[yǎnchū]{yǎnchū — to perform on stage}

\label{lesson:ZH-C257-yanchu}



\begin{quote}
Before the new one: say the Mandarin for a musical instrument, then the Mandarin for to play the drums.
\end{quote}

\subsection*{You'll want to know: yǎnchū}
\emph{yǎnchū} — \textquotedblleft{}to perform on stage\textquotedblright{}.

\begin{quote}
\emph{jīntiān wǎnshang tāmen yǎnchū} — They're performing tonight.
\end{quote}

Its characters come later; for now it is a word for the ear.

\begin{cousinweb}
\emph{yǎnchū} is built from \emph{yǎn} and \zh{出} 'to come out', both inherited from Middle Chinese, \emph{yenX} and \emph{tsyhwit} (Guangyun, in Baxter's transcription), which glosses the first 'wide' and 'water flowing far'. The sense 'to perform' is not given there, so how it came is not settled.
\end{cousinweb}

\begin{grammarlens}[title={What you've built}]
One more word for what I like.
\end{grammarlens}

\subsection*{Your turn}
\begin{itemize}
\raggedright
  \item \emph{Say it:} \emph{yǎnchū}
  \item \emph{Say it:} \emph{yǎnchū}, once more
  \item \emph{Say it:} \emph{jīntiān wǎnshang tāmen yǎnchū} — They're performing tonight.
  \item \emph{Recall:} say the Mandarin for a musical instrument, then the Mandarin for to play the drums, then say \emph{yǎnchū} again
\end{itemize}

\begin{tcolorbox}[breakable,colback=teal!4,colframe=teal!35!black,title={Before you move on}]
What does yǎnchū mean? (\textquotedblleft{}To perform on stage\textquotedblright{}.) Say it once more.
\end{tcolorbox}
\section[gēshǒu]{gēshǒu — singer}

\label{lesson:ZH-C257-geshou}



\begin{quote}
Before the new one: say the Mandarin for to play the drums, then the Mandarin for to perform on stage.
\end{quote}

\subsection*{You'll want to know: gēshǒu}
\emph{gēshǒu} — \textquotedblleft{}singer\textquotedblright{}.

\begin{quote}
\emph{tā shì wǒ zuì xǐhuan de gēshǒu} — He's my favourite singer.
\end{quote}

Its characters come later; for now it is a word for the ear.

\begin{cousinweb}
\emph{gēshǒu} is built from \emph{gē} 'a song' and \zh{手} 'a hand', both inherited from Middle Chinese, \emph{ka} and \emph{syuwX} (Guangyun, in Baxter's transcription), which quotes an older book for the first: the voice of people is called song. A hand here is someone skilled at a thing; that sense is not given there.
\end{cousinweb}

\begin{grammarlens}[title={What you've built}]
One more word for what I like.
\end{grammarlens}

\subsection*{Your turn}
\begin{itemize}
\raggedright
  \item \emph{Say it:} \emph{gēshǒu}
  \item \emph{Say it:} \emph{gēshǒu}, once more
  \item \emph{Say it:} \emph{tā shì wǒ zuì xǐhuan de gēshǒu} — He's my favourite singer.
  \item \emph{Recall:} say the Mandarin for to play the drums, then the Mandarin for to perform on stage, then say \emph{gēshǒu} again
\end{itemize}

\begin{tcolorbox}[breakable,colback=teal!4,colframe=teal!35!black,title={Before you move on}]
What does gēshǒu mean? (\textquotedblleft{}Singer\textquotedblright{}.) Say it once more.
\end{tcolorbox}
\section[gǔzhǎng]{gǔzhǎng — to clap; to applaud}

\label{lesson:ZH-C257-guzhang}



\begin{quote}
Before the new one: say the Mandarin for to perform on stage, then the Mandarin for a singer.
\end{quote}

\subsection*{You'll want to know: gǔzhǎng}
\emph{gǔzhǎng} — \textquotedblleft{}to clap; to applaud\textquotedblright{}.

\begin{quote}
\emph{dàjiā dōu gǔzhǎng le} — Everyone clapped.
\end{quote}

Its characters come later; for now it is a word for the ear.

Two third tones meet in this word: gǔ is said with a rising tone before zhǎng, and the spelling keeps both marks.

\begin{cousinweb}
\emph{gǔzhǎng} is built from \emph{gǔ} 'a drum' and \emph{zhǎng} 'the palm of the hand', both inherited from Middle Chinese, \emph{kuX} and \emph{tsyangX} (Guangyun, in Baxter's transcription), which quotes an older dictionary for the first, 'to beat a drum', and glosses the second 'the palm'.
\end{cousinweb}

\begin{grammarlens}[title={What you've built}]
One more word for what I like.
\end{grammarlens}

\subsection*{Your turn}
\begin{itemize}
\raggedright
  \item \emph{Say it:} \emph{gǔzhǎng}
  \item \emph{Say it:} \emph{gǔzhǎng}, once more
  \item \emph{Say it:} \emph{dàjiā dōu gǔzhǎng le} — Everyone clapped.
  \item \emph{Recall:} say the Mandarin for to perform on stage, then the Mandarin for a singer, then say \emph{gǔzhǎng} again
\end{itemize}

\begin{tcolorbox}[breakable,colback=teal!4,colframe=teal!35!black,title={Before you move on}]
What does gǔzhǎng mean? (\textquotedblleft{}To clap; to applaud\textquotedblright{}.) Say it once more.
\end{tcolorbox}
